% !TEX root = ../main.tex
\begin{abstract}
\noindent Bantay-Baha ranks Philippine provinces for action before a typhoon arrives. It combines the probability $R$ that a province's dengue cases surge in the next four weeks with the share $A$ of its people that the storm cuts off from a hospital, and ranks provinces by the priority $P = R \times A$. This specification describes the data, the method of every layer and how each was evaluated; every result in it is computed from the pipeline outputs at build time. Over \nScored{} rolling-origin forecasts in 2018--2019, a multi-quantile gradient-boosted model reaches a relative weighted interval score (WIS) of \relWIS{} against a seasonal-naive forecast, with a 90\% interval coverage of \covNinety{}; its median ensemble with \nBench{} benchmark models reaches \ensRelWIS{}. For the four backtest typhoons, \cutNetRange{} people lose their route to a hospital on the road and ferry network when roads are cut by flood observed from radar. The lead-lag outbreak graph between provinces keeps \llEdges{} links; at one to four weeks its features \ifGraphHelps improve\else do not improve\fi{} the forecast, and the rain and flood features \ifEnvHelps do\else do not either\fi.
\end{abstract}
